\documentclass{article}

% Packages
\usepackage{microtype}
\usepackage{enumitem}
\usepackage{pgffor}

% Package stuff
\setlist{nosep}

% Document attributes
\title{
    Game Design Document\\
    Nevermind
}
\author{Andrea Ricciardi \and Christian Vezzoli \and Shaan Vashisht \and Matteo Vannini \and Andrea Coscia Vanoni}

% Document
\begin{document}

\maketitle
\clearpage

\tableofcontents
\clearpage

\section{Overview and vision statement}

\section{Gameplay}
\subsection{Main Mechanics}
\subsubsection{Player Mechanics}
\begin{itemize}
    \item \textbf{Movement:} The player has a physics based character controller that allows them to move sideways.
    \item \textbf{Memory System:} The player has a set of memories that they can remember/forget. 
    The elements in the level change their behavior depending on the memories the player currently remembers.
    When the player forgets a memory, a thought bubble containing it is dropped at their position. The player has to interact with the bubble to remember the memory.
\end{itemize}

\subsubsection{Level Elements}
\begin{itemize}
    \item \textbf{Pressure Plate:} Sends signal when one or more objects are placed on top (the amount of weight needed can be specified).
    \item \textbf{Button:} Sends signal when the player collides with it (only once).
    \item \textbf{Lever:} Sends signal when the player collides with it (can switch on/off repeatedly).
    \item \textbf{Moving Platform:} Platform that moves along a set path.
    \item \textbf{Rotating Platform:} Platform that rotates.
    \item \textbf{Gate:} Opens/Closes when signaled.
    \item \textbf{Box:} Block that can be pushed or grabbed by the player.
    \item \textbf{Ball:} Can be picked up and thrown.
    \item \textbf{Kill Zone:} Kills the player on contact.
    \item \textbf{Slime*:} Surface mod, makes objects stick together. Big slime area can be swam through.
    \item \textbf{Light Beam*:} Emits a ray of light, can be rotated by the player.
    \item \textbf{Light Beam* - Solid:} Ray of light that acts as a physical wall/platform.
    \item \textbf{Light Beam* - Deadly:} Ray of light that kills the player.
    \item \textbf{Light Receiver:} Sends signal when reached by a ray of light.
    \item \textbf{Mirror:} Redirects light.
    \item \textbf{Memory Screen:} Makes the player forget fragile memories when passed through.
    \item \textbf{Memory Font:} Restores the player's semi-permanent memories.
    \item \textbf{Ice*:} Surface mod, makes objects slide.
    \item \textbf{Rubber*:} Surface mod, makes objects bounce.
    \item \textbf{Rocket:} Kills player if hit, can be used as a platform.
    \item \textbf{Gravity* Flux:} Pushes objects in a given direction.
    \item \textbf{Gravity* Point:} Attracts objects towards the point.
\end{itemize}
\textit{*tied to a memory that enables/disable the object.}

\subsubsection{Memories}
\begin{itemize}
    \item \textbf{Material (eg. Stone):} Manifests objects of a given material.
    \item \textbf{Movement:} Allows the player to move.
    \item \textbf{Jumping:} Allows the player to jump/Makes the player jump uncontrollably.
    \item \textbf{Time:} Pauses the game when forgotten.
    \item \textbf{Hands:} Allows the player to pick up and throw small objects (eg. balls, small boxes, memories).
    \item \textbf{Object Specific:} See previous section. 
\end{itemize}

\paragraph{Memory Types}
\begin{itemize}
    \item \textbf{Permanent:} Cannot be forgotten.
    \item \textbf{Semi-Permanent:} Can be forgotten, will return to the player at a memory font.
    \item \textbf{Fragile:} Can be forgotten, will be destroyed when passing through a memory screen.
\end{itemize}

\subsection{Level Design}
A level is defined as a single scene composed of a succession of rooms, each occupying one screen and containing one puzzle.
Each level is designed around a theme (eg. "forest", "volcano", "temple") and a central mechanic (eg. "box-pushing", "slime", "lights") that is explored throughout the level's rooms.

\section{Characters*}

\section{Story*}

\section{World*}

\section{Media list}

\section{Technical specification*}

\section{Team}
\begin{itemize}
    \item Andrea Ruben Ricciardi
    \item Christian Vezzoli
    \item Shaan Edoardo Vashisht
    \item Matteo Vannini
    \item Andrea Coscia Vanoni
\end{itemize}

\section{Special Thanks}

\section{Team meetings and Deadlines}
\foreach \i in {1, 2, ..., 20} {%
    \edef\FileName{Meeting Notes/Week\i}%
    \IfFileExists{\FileName}{%
       \input{\FileName}%
    }
}

\end{document}